Figure~\ref{fig:schematic} summarises the study. This study asks a question that the literature has assumed rather than tested: does making a flight policy tail-risk sensitive make a multi-UAV delivery operation more resilient to distribution shift, and if not, what does? We answer it with a controlled, paired experiment. We build a stochastic urban delivery simulator in which a fleet delivers parcels under wind, turbulence and gusts, satellite-navigation drift and multipath, cooperative and non-cooperative traffic, and uncertain battery energy, payload and propulsion faults, and in which four shift families can be dialled independently at test time. Within this testbed we train six learning agents under a matched budget: the risk-neutral soft actor-critic, TD3 and PPO, a risk-neutral distributional agent, a fixed-risk CVaR agent of the kind used in risk-sensitive navigation, and our hazard-adaptive agent CREST, which conditions both actor and critics on a risk level that is scheduled online from four hazard cues, namely collision threat, energy margin, traffic density and disturbance intensity. We study two placements of the tail-risk measure, on the distribution of total returns as in previous work and on the distribution of a catastrophe cost, and we train the risk-sensitive agents with a soft-risk curriculum and with risk-level relabelling so that a single policy covers a continuum of risk preferences.

Alongside the policy we treat the return-to-home decision as a statistical problem rather than a fixed rule of thumb. We introduce an online conformal energy reserve in which two per-aircraft quantile trackers follow the ratio of measured to modelled power and the along-track wind innovation, and a fleet-level calibrator adjusts the reserve once per completed return leg so that the long-run fraction of legs that exceed their predicted energy matches a target. The construction inherits the distribution-free long-run guarantee of adaptive conformal inference, which we state and prove for the leg-level calibrator, and it requires no assumption that test conditions resemble training conditions.

The evaluation is paired at the level of individual scenarios, uses three training seeds and about one thousand missions per policy and condition, and separates policy effects from reserve effects by running every policy under every reserve rule. Our findings are not the ones the risk-sensitive navigation literature would predict. At a matched training budget, none of the three tail-risk designs delivers a reliable reduction in catastrophic losses relative to a well-tuned risk-neutral soft actor-critic: the gaps stay inside the bootstrap intervals, change sign across conditions and are of the same size as the spread across training seeds. Turning the risk dial towards caution leaves collision rates unchanged or slightly worse, because caution expresses itself as hesitancy that lengthens exposure to conflicts rather than as wider separation. The decisive lever is instead the energy reserve: the conformal reserve attains its specified coverage in every shift family whereas fixed margins silently lose coverage as conditions drift, and at operating points with negligible depletion it completes more deliveries than the fixed margin that is equally safe. We also show where the conformal reserve is not the better choice, quantify the abort-depletion frontier that an operator must choose on, and identify conflicts with other aircraft under traffic and sensing shift as the residual failure mode that neither risk-sensitive learning nor energy calibration addresses. Table~\ref{tab:compare} places the study relative to representative risk-sensitive and calibration-based approaches. The Methods describe the environment, the agents, the reserves and the evaluation protocol, the Results and Discussion follow, and the Conclusion summarises the implications. Taken together, the results argue for putting statistical calibration where the failure mechanism is understood, and for treating claims of resilience from tail-risk objectives as hypotheses to be tested against strong risk-neutral baselines.
